% =============================================================
% Practice Problems: Instrumental Variables Regression
% Course: Quantitative Economics (PhD), Vilnius University
% Author: Swapnil Singh
% =============================================================
\documentclass[10pt]{beamer}

\usepackage[T1]{fontenc}
\usepackage{graphicx}
\usepackage{booktabs}
\usepackage{tabularx}
\usepackage{array}
\usepackage{ragged2e}
\usepackage{appendixnumberbeamer}
\usepackage{amsmath,amssymb,amsfonts}
\usepackage{mathtools}
\usepackage{cancel}
\usepackage{hyperref}
\usepackage{pifont}
\usepackage{tikz}
\usetikzlibrary{arrows.meta,positioning,calc}

% SimplePlus theme
\usetheme{SimplePlus}

\title{Practice Problems: Instrumental Variables}
\subtitle{Quantitative Economics}
\author{Swapnil Singh}
\institute{Lietuvos Bankas \and Kaunas University of Technology}
\date{Vilnius University, 2026}

\begin{document}

% ---- TITLE SLIDE ----
\begin{frame}[plain]
  \titlepage
  \vfill
  {\scriptsize\color{gray} 20 problems from simple to tough. Solutions in the appendix.}
\end{frame}

% ---- TABLE OF CONTENTS ----
\begin{frame}{Overview}
\textbf{20 problems}, ordered by difficulty:
\begin{itemize}
  \item Problems 1--5: Definitions \& conceptual foundations
  \item Problems 6--10: Bias formulas \& derivations
  \item Problems 11--15: TSLS mechanics \& consistency
  \item Problems 16--20: Advanced topics \& applications
\end{itemize}

\vspace{1em}
Each problem slide has a \beamerbutton{Solution} button linking to the appendix.\\
Each solution slide has a \beamerbutton{Back} button to return.
\end{frame}

% ============================================================
%  QUESTIONS
% ============================================================

% ---- Q1 ----
\begin{frame}{Problem 1: Three sources of endogeneity}
\label{q1}
List the three classical sources of endogeneity. For each, give a one-sentence explanation of \emph{why} it causes $\operatorname{cov}(X_i, u_i) \neq 0$.

\vfill
\hfill\hyperlink{s1}{\beamerbutton{Solution}}
\end{frame}

% ---- Q2 ----
\begin{frame}{Problem 2: Instrument validity conditions}
\label{q2}
State the two conditions required for a variable $Z$ to be a valid instrument for the endogenous regressor $X$ in the model $Y_i = \beta_0 + \beta_1 X_i + u_i$.

For each condition, explain in one sentence what goes wrong if it fails.

\vfill
\hfill\hyperlink{s2}{\beamerbutton{Solution}}
\end{frame}

% ---- Q3 ----
\begin{frame}{Problem 3: Endogenous or exogenous?}
\label{q3}
For each scenario, state whether $X$ is endogenous or exogenous and identify which source of endogeneity (if any) applies:

\begin{enumerate}[(a)]
  \item Regressing wages on years of education, omitting innate ability.
  \item Regressing crop yield on fertilizer use, where fertilizer is randomly assigned across plots.
  \item Regressing quantity demanded on price in a competitive market.
  \item Regressing test scores on hours studied, where hours studied is measured with error.
\end{enumerate}

\vfill
\hfill\hyperlink{s3}{\beamerbutton{Solution}}
\end{frame}

% ---- Q4 ----
\begin{frame}{Problem 4: Levels of exogeneity}
\label{q4}
State the three progressively stronger conditions for instrument exogeneity:
\begin{enumerate}[(a)]
  \item Zero correlation
  \item Mean independence
  \item Full independence
\end{enumerate}

\vspace{0.5em}
Show that mean independence implies zero correlation. Then provide a counterexample demonstrating that the converse is false.

\vfill
\hfill\hyperlink{s4}{\beamerbutton{Solution}}
\end{frame}

% ---- Q5 ----
\begin{frame}{Problem 5: Evaluating proposed instruments}
\label{q5}
A researcher wants to estimate the effect of class size on student test scores. She proposes the following instruments. For each, discuss whether relevance and exogeneity are likely satisfied:

\begin{enumerate}[(a)]
  \item School district budget per pupil.
  \item Random fluctuations in cohort size due to birth-date variation.
  \item Average temperature in the school district.
\end{enumerate}

\vfill
\hfill\hyperlink{s5}{\beamerbutton{Solution}}
\end{frame}

% ---- Q6 ----
\begin{frame}{Problem 6: Omitted variable bias formula}
\label{q6}
Suppose the true DGP is:
\[
Y_i = \beta_0 + \beta_1 X_i + \beta_2 W_i + \varepsilon_i, \qquad E(\varepsilon_i \mid X_i, W_i) = 0
\]
but $W_i$ is unobserved. Define $u_i = \beta_2 W_i + \varepsilon_i$.

\begin{enumerate}[(a)]
  \item Derive $\operatorname{cov}(X_i, u_i)$ and state conditions under which it is nonzero.
  \item Show that $\widehat{\beta}_1^{OLS} \xrightarrow{p} \beta_1 + \beta_2 \dfrac{\operatorname{cov}(X_i, W_i)}{\operatorname{var}(X_i)}$.
  \item If $\beta_2 > 0$ and $\operatorname{cov}(X,W) > 0$, is the OLS estimate biased upward or downward?
\end{enumerate}

\vfill
\hfill\hyperlink{s6}{\beamerbutton{Solution}}
\end{frame}

% ---- Q7 ----
\begin{frame}{Problem 7: Attenuation bias}
\label{q7}
True model: $Y_i = \beta_0 + \beta_1 X_i^* + \varepsilon_i$ with $E(\varepsilon_i \mid X_i^*) = 0$.

We observe $X_i = X_i^* + w_i$ (classical measurement error).

\begin{enumerate}[(a)]
  \item Substitute $X_i^* = X_i - w_i$ and derive the composite error $u_i$.
  \item Show that $\operatorname{cov}(X_i, u_i) = -\beta_1 \sigma_w^2$.
  \item Derive the probability limit of $\widehat{\beta}_1^{OLS}$ and explain why the bias is called ``attenuation bias.''
  \item If the true $\beta_1 = 2$, $\sigma_{X^*}^2 = 4$, and $\sigma_w^2 = 1$, compute the plim of OLS.
\end{enumerate}

\vfill
\hfill\hyperlink{s7}{\beamerbutton{Solution}}
\end{frame}

% ---- Q8 ----
\begin{frame}{Problem 8: Simultaneous equations bias}
\label{q8}
Consider the supply--demand system:
\begin{align*}
\text{Demand}: \quad Q_i &= \alpha_0 + \alpha_1 P_i + u_i^d \qquad (\alpha_1 < 0) \\
\text{Supply}: \quad Q_i &= \gamma_0 + \gamma_1 P_i + u_i^s \qquad (\gamma_1 > 0)
\end{align*}
Assume $u^d \perp u^s$, both mean-zero with variances $\sigma_d^2$ and $\sigma_s^2$.

\begin{enumerate}[(a)]
  \item Solve for the equilibrium price $P_i$.
  \item Show that $\operatorname{cov}(P_i, u_i^d) = \dfrac{-\sigma_d^2}{\alpha_1 - \gamma_1}$ and explain its sign.
  \item Derive: $\widehat{\alpha}_1^{OLS} \xrightarrow{p} \dfrac{\alpha_1 \sigma_s^2 + \gamma_1 \sigma_d^2}{\sigma_s^2 + \sigma_d^2}$.
  \item Interpret this result: what does OLS estimate?
\end{enumerate}

\vfill
\hfill\hyperlink{s8}{\beamerbutton{Solution}}
\end{frame}

% ---- Q9 ----
\begin{frame}{Problem 9: Exogeneity failure and TSLS inconsistency}
\label{q9}
Suppose the instrument $Z$ is not perfectly exogenous: $\operatorname{cov}(Z_i, u_i) = \rho \neq 0$.

\begin{enumerate}[(a)]
  \item Show that $\widehat{\beta}_1^{TSLS} \xrightarrow{p} \beta_1 + \dfrac{\rho}{\operatorname{cov}(Z_i, X_i)}$.
  \item If $\rho = 0.01$ and $\operatorname{cov}(Z_i, X_i) = 0.5$, what is the asymptotic bias?
  \item Now suppose the instrument is weak: $\operatorname{cov}(Z_i, X_i) = 0.02$. Recompute the asymptotic bias. What do you observe?
  \item Explain in words why weak instruments \emph{amplify} violations of exogeneity.
\end{enumerate}

\vfill
\hfill\hyperlink{s9}{\beamerbutton{Solution}}
\end{frame}

% ---- Q10 ----
\begin{frame}{Problem 10: Direction of OVB with IV}
\label{q10}
Return to the OVB setup. The true DGP is $Y_i = \beta_0 + \beta_1 X_i + \beta_2 W_i + \varepsilon_i$.

An instrument $Z$ satisfies $\operatorname{cov}(Z,X) \neq 0$ and $\operatorname{cov}(Z, \varepsilon) = 0$.

\begin{enumerate}[(a)]
  \item Write the exogeneity condition $\operatorname{cov}(Z_i, u_i) = 0$ in terms of $\beta_2$ and $\operatorname{cov}(Z_i, W_i)$.
  \item Why does this mean the instrument must be uncorrelated with the omitted variable $W$?
  \item A researcher proposes parental income as an instrument for education in a wage regression (omitting ability). Is this a good instrument? Why or why not?
\end{enumerate}

\vfill
\hfill\hyperlink{s10}{\beamerbutton{Solution}}
\end{frame}

% ---- Q11 ----
\begin{frame}{Problem 11: The TSLS formula}
\label{q11}
In the simple model $Y_i = \beta_0 + \beta_1 X_i + u_i$ with one instrument $Z$:

\begin{enumerate}[(a)]
  \item The TSLS estimator is the OLS estimator of $Y$ on $\widehat{X}$, where $\widehat{X}_i = \widehat{\pi}_0 + \widehat{\pi}_1 Z_i$. Starting from $\widehat{\beta}_1^{TSLS} = \dfrac{s_{\widehat{X}Y}}{s_{\widehat{X}}^2}$, show that $\widehat{\beta}_1^{TSLS} = \dfrac{s_{ZY}}{s_{ZX}}$.

  \item Explain why this is called the ``Wald estimator'' and interpret it as a ratio of a reduced-form effect to a first-stage effect.
\end{enumerate}

\vfill
\hfill\hyperlink{s11}{\beamerbutton{Solution}}
\end{frame}

% ---- Q12 ----
\begin{frame}{Problem 12: Consistency of TSLS}
\label{q12}
Starting from $Y_i = \beta_0 + \beta_1 X_i + u_i$:

\begin{enumerate}[(a)]
  \item Take population covariances of both sides with $Z_i$ and solve for $\beta_1$ under the IV assumptions.
  \item Use the law of large numbers and continuous mapping theorem to show that $\widehat{\beta}_1^{TSLS} = s_{ZY}/s_{ZX} \xrightarrow{p} \beta_1$.
  \item At which step does instrument relevance play a role? At which step does instrument exogeneity play a role?
\end{enumerate}

\vfill
\hfill\hyperlink{s12}{\beamerbutton{Solution}}
\end{frame}

% ---- Q13 ----
\begin{frame}{Problem 13: First-stage $F$-statistic}
\label{q13}
A researcher estimates a first-stage regression and obtains a first-stage $F$-statistic of 7.2.

\begin{enumerate}[(a)]
  \item Is the instrument weak according to the rule of thumb?
  \item Using the approximate bias formula $E(\widehat{\beta}_1^{TSLS}) - \beta_1 \approx \dfrac{\beta_1^{OLS} - \beta_1}{E(F) - 1}$, what fraction of the OLS bias does the TSLS estimate retain?
  \item List three strategies the researcher could pursue.
  \item Why are standard TSLS confidence intervals unreliable with weak instruments?
\end{enumerate}

\vfill
\hfill\hyperlink{s13}{\beamerbutton{Solution}}
\end{frame}

% ---- Q14 ----
\begin{frame}{Problem 14: Identification and the order condition}
\label{q14}
Consider the general IV model:
\[
Y_i = \beta_0 + \beta_1 X_{1i} + \beta_2 X_{2i} + \beta_3 W_{1i} + u_i
\]
where $X_1, X_2$ are endogenous and $W_1$ is exogenous.

\begin{enumerate}[(a)]
  \item How many instruments are needed at minimum?
  \item A researcher proposes three instruments $Z_1, Z_2, Z_3$. Is the model exactly identified, overidentified, or underidentified?
  \item Describe what the first-stage regressions look like. How many are there?
  \item Can the researcher test instrument exogeneity? If so, what is the degrees of freedom of the $J$-test?
\end{enumerate}

\vfill
\hfill\hyperlink{s14}{\beamerbutton{Solution}}
\end{frame}

% ---- Q15 ----
\begin{frame}{Problem 15: The $J$-test}
\label{q15}
A researcher estimates a model with $k=1$ endogenous regressor and $m=3$ instruments. After estimating by TSLS:

\begin{enumerate}[(a)]
  \item Describe the steps to compute the $J$-statistic.
  \item Under $H_0$, what is the distribution of the $J$-statistic? State the degrees of freedom.
  \item The researcher obtains $J = 7.8$. At the 5\% significance level, does she reject? (Use $\chi^2_2(0.05) = 5.99$.)
  \item The $J$-test rejects. Does this mean \emph{all} instruments are endogenous? Explain.
\end{enumerate}

\vfill
\hfill\hyperlink{s15}{\beamerbutton{Solution}}
\end{frame}

% ---- Q16 ----
\begin{frame}{Problem 16: Sampling distribution of TSLS}
\label{q16}
From the representation $\widehat{\beta}_1^{TSLS} = \beta_1 + \dfrac{\frac{1}{n}\sum(Z_i - \bar{Z})u_i}{\frac{1}{n}\sum(Z_i - \bar{Z})(X_i - \bar{X})}$:

\begin{enumerate}[(a)]
  \item Define $\bar{q} = \frac{1}{n}\sum(Z_i - \bar{Z})u_i$. Show that $E(\bar{q}) = 0$ under instrument exogeneity and state which theorem implies $\sqrt{n}\bar{q} \xrightarrow{d} N(0, \sigma_q^2)$.
  \item The denominator converges to $\operatorname{cov}(Z,X)$ in probability. Which theorem combines numerator and denominator to yield asymptotic normality of $\widehat{\beta}_1^{TSLS}$?
  \item Write down the asymptotic variance $\sigma_{\widehat{\beta}_1^{TSLS}}^2$ and explain why it is larger when $\operatorname{cov}(Z,X)$ is small.
\end{enumerate}

\vfill
\hfill\hyperlink{s16}{\beamerbutton{Solution}}
\end{frame}

% ---- Q17 ----
\begin{frame}{Problem 17: Cigarette demand application}
\label{q17}
Using panel data (10-year differences), a researcher estimates:
\[
\Delta\ln Q_i = \beta_0 + \beta_1 \Delta\ln P_i + \beta_2 \Delta\ln\text{Inc}_i + u_i
\]
with instruments $\Delta\text{SalesTax}$ and $\Delta\text{CigTax}$.

Results: $\widehat{\beta}_1^{(1)} = -0.94$ (SalesTax only), $\widehat{\beta}_1^{(2)} = -1.34$ (CigTax only), $\widehat{\beta}_1^{(3)} = -1.20$ (both). The $J$-test $p$-value with both instruments is 0.026.

\begin{enumerate}[(a)]
  \item First-stage $F$-statistics are 33.7, 107.2, and 88.6. Are instruments weak?
  \item Interpret the $J$-test result. What does rejection imply?
  \item Which instrument would you prefer and why? What is the economic argument?
  \item Interpret the preferred elasticity estimate in terms of policy.
\end{enumerate}

\vfill
\hfill\hyperlink{s17}{\beamerbutton{Solution}}
\end{frame}

% ---- Q18 ----
\begin{frame}{Problem 18: Anderson--Rubin test}
\label{q18}
The Anderson--Rubin (AR) test of $H_0: \beta_1 = \beta_{1,0}$ works by:
\begin{enumerate}
  \item Constructing $Y_i^* = Y_i - \beta_{1,0}X_i$
  \item Regressing $Y_i^*$ on $Z$'s and $W$'s
  \item Testing whether all coefficients on $Z$'s are jointly zero
\end{enumerate}

\begin{enumerate}[(a)]
  \item Explain intuitively why, if $\beta_1 = \beta_{1,0}$ and the instruments are exogenous, the $Z$'s should have no predictive power for $Y_i^*$.
  \item Why is this test valid even when instruments are weak?
  \item Describe how to construct a confidence set by inverting the AR test.
  \item What unusual features can AR confidence sets have?
\end{enumerate}

\vfill
\hfill\hyperlink{s18}{\beamerbutton{Solution}}
\end{frame}

% ---- Q19 ----
\begin{frame}{Problem 19: IV with simultaneous causality -- full derivation}
\label{q19}
In the supply--demand system with a supply shifter $Z$:
\begin{align*}
\text{Demand}: & \quad Q_i = \alpha_0 + \alpha_1 P_i + u_i^d \\
\text{Supply}: & \quad Q_i = \gamma_0 + \gamma_1 P_i + \delta Z_i + u_i^s
\end{align*}
Assume $\operatorname{cov}(Z_i, u_i^d) = 0$, $\operatorname{cov}(Z_i, u_i^s) = 0$, and $\delta \neq 0$.

\begin{enumerate}[(a)]
  \item Solve for the equilibrium price $P_i$ in terms of $Z_i$ and the shocks.
  \item Verify that $\operatorname{cov}(Z_i, P_i) \neq 0$ (relevance).
  \item Verify that $\operatorname{cov}(Z_i, u_i^d) = 0$ (exogeneity for the demand equation).
  \item Compute $\operatorname{cov}(Z_i, Q_i)$ using the demand equation and show that $\dfrac{\operatorname{cov}(Z_i, Q_i)}{\operatorname{cov}(Z_i, P_i)} = \alpha_1$, confirming TSLS identifies the demand slope.
\end{enumerate}

\vfill
\hfill\hyperlink{s19}{\beamerbutton{Solution}}
\end{frame}

% ---- Q20 ----
\begin{frame}{Problem 20: Comprehensive analysis}
\label{q20}
A health economist estimates the effect of hospital spending ($X$) on patient survival ($Y$). She suspects spending is endogenous because sicker patients receive more spending.

She proposes two instruments: (i) distance to nearest teaching hospital ($Z_1$), (ii) local property tax revenue ($Z_2$). Her results:
\begin{itemize}
  \item OLS: $\widehat{\beta}_1^{OLS} = -0.05$ (SE $= 0.01$)
  \item TSLS (both IVs): $\widehat{\beta}_1^{TSLS} = 0.12$ (SE $= 0.04$)
  \item First-stage $F = 14.3$
  \item $J$-statistic $= 0.45$, $p$-value $= 0.50$
\end{itemize}

\begin{enumerate}[(a)]
  \item Why is OLS likely biased? In which direction?
  \item Assess instrument strength. Are the instruments weak?
  \item Interpret the $J$-test. Does it validate both instruments?
  \item Discuss the exogeneity of each instrument. Which is more credible?
  \item Reconcile the OLS and TSLS estimates economically.
\end{enumerate}

\vfill
\hfill\hyperlink{s20}{\beamerbutton{Solution}}
\end{frame}

% ============================================================
%  APPENDIX: SOLUTIONS
% ============================================================
\appendix

\begin{frame}[plain]
  \centering
  \vfill
  {\Huge\bfseries Solutions}
  \vfill
\end{frame}

% ---- S1 ----
\begin{frame}{Solution 1: Three sources of endogeneity}
\label{s1}
\begin{enumerate}
  \item \textbf{Omitted Variable Bias}: An unobserved variable $W$ in the error $u$ is correlated with $X$, so $\operatorname{cov}(X, u) = \beta_2\operatorname{cov}(X, W) \neq 0$.

  \item \textbf{Measurement Error}: The observed $X = X^* + w$ contains noise $w$ that mechanically enters both $X$ and $u = \varepsilon - \beta_1 w$, creating $\operatorname{cov}(X, u) = -\beta_1\sigma_w^2 \neq 0$.

  \item \textbf{Simultaneous Causality}: $X$ and $Y$ are jointly determined in equilibrium, so the structural error $u$ feeds back into $X$ through the equilibrium mapping.
\end{enumerate}

\vfill
\hfill\hyperlink{q1}{\beamerbutton{Back to Problem 1}}
\end{frame}

% ---- S2 ----
\begin{frame}{Solution 2: Instrument validity conditions}
\label{s2}
\begin{block}{Condition 1: Instrument Relevance}
$\operatorname{corr}(Z_i, X_i) \neq 0$. \quad $Z$ must be correlated with $X$.

\emph{If it fails}: $Z$ carries no information about $X$. The denominator $\operatorname{cov}(Z,X) \to 0$, making the TSLS estimator undefined (ratio of two normals, not consistent).
\end{block}

\begin{block}{Condition 2: Instrument Exogeneity}
$\operatorname{corr}(Z_i, u_i) = 0$. \quad $Z$ must be uncorrelated with the error.

\emph{If it fails}: TSLS converges to $\beta_1 + \operatorname{cov}(Z,u)/\operatorname{cov}(Z,X) \neq \beta_1$. The estimator is inconsistent --- it ``locks onto'' contaminated variation.
\end{block}

\vfill
\hfill\hyperlink{q2}{\beamerbutton{Back to Problem 2}}
\end{frame}

% ---- S3 ----
\begin{frame}{Solution 3: Endogenous or exogenous?}
\label{s3}
\begin{enumerate}[(a)]
  \item \textbf{Endogenous} --- \emph{Omitted variable bias}. Ability is in $u$, and $\operatorname{cov}(\text{education}, \text{ability}) > 0$.

  \item \textbf{Exogenous}. Random assignment ensures $\operatorname{cov}(X, u) = 0$. OLS is consistent.

  \item \textbf{Endogenous} --- \emph{Simultaneous causality}. Price and quantity are jointly determined by supply and demand; OLS on either equation is inconsistent.

  \item \textbf{Endogenous} --- \emph{Measurement error}. Observed $X = X^* + w$, and the noise $w$ enters both $X$ and the composite error, creating $\operatorname{cov}(X, u) = -\beta_1\sigma_w^2 \neq 0$.
\end{enumerate}

\vfill
\hfill\hyperlink{q3}{\beamerbutton{Back to Problem 3}}
\end{frame}

% ---- S4 ----
\begin{frame}{Solution 4: Levels of exogeneity}
\label{s4}
\textbf{Three levels}: (a) $\operatorname{cov}(Z,u)=0$; \quad (b) $E(u\mid Z)=0$; \quad (c) $u \perp\!\!\!\perp Z$.

\textbf{Proof that (b) $\Rightarrow$ (a)}:
\begin{align*}
E(Zu) &= E[E(Zu \mid Z)] = E[Z \cdot \underbrace{E(u \mid Z)}_{=0}] = 0 \\
E(u) &= E[\underbrace{E(u \mid Z)}_{=0}] = 0
\end{align*}
So $\operatorname{cov}(Z,u) = E(Zu) - E(Z)E(u) = 0$.

\textbf{Counterexample (converse fails)}: Let $Z \sim N(0,1)$, $u = Z^2 - 1$.
\begin{itemize}
  \item $E(u) = 0$, $\operatorname{cov}(Z,u) = E(Z^3) = 0$. \checkmark
  \item But $E(u \mid Z = z) = z^2 - 1 \neq 0$ for $z \neq \pm 1$. \ding{55}
\end{itemize}

\vfill
\hfill\hyperlink{q4}{\beamerbutton{Back to Problem 4}}
\end{frame}

% ---- S5 ----
\begin{frame}{Solution 5: Evaluating proposed instruments}
\label{s5}
\begin{enumerate}[(a)]
  \item \textbf{Budget per pupil}: \emph{Relevant} (more funding $\to$ smaller classes), but likely \emph{not exogenous} --- budget correlates with district wealth, parental involvement, and school quality, all of which affect test scores directly.

  \item \textbf{Random birth-date fluctuations} (Hoxby, 2000): \emph{Relevant} (more students $\to$ larger classes) and plausibly \emph{exogenous} --- biological variation in birth dates is unrelated to school quality or family background. \alert{Good instrument.}

  \item \textbf{Average temperature}: Likely \emph{irrelevant} (no clear mechanism linking temperature to class size). Even if correlated, the link is too weak. Not a useful instrument.
\end{enumerate}

\vfill
\hfill\hyperlink{q5}{\beamerbutton{Back to Problem 5}}
\end{frame}

% ---- S6 ----
\begin{frame}{Solution 6: Omitted variable bias formula}
\label{s6}
\begin{enumerate}[(a)]
  \item $\operatorname{cov}(X, u) = \operatorname{cov}(X, \beta_2 W + \varepsilon) = \beta_2\operatorname{cov}(X,W) + \underbrace{\operatorname{cov}(X,\varepsilon)}_{=0} = \beta_2\operatorname{cov}(X,W)$.

  Nonzero iff $\beta_2 \neq 0$ \emph{and} $\operatorname{cov}(X,W) \neq 0$.

  \item $\widehat{\beta}_1^{OLS} \xrightarrow{p} \beta_1 + \dfrac{\operatorname{cov}(X,u)}{\operatorname{var}(X)} = \beta_1 + \beta_2\dfrac{\operatorname{cov}(X,W)}{\operatorname{var}(X)}$. \checkmark

  \item $\beta_2 > 0$ and $\operatorname{cov}(X,W) > 0$ $\Rightarrow$ bias $= \beta_2\dfrac{\operatorname{cov}(X,W)}{\operatorname{var}(X)} > 0$.

  OLS is \alert{biased upward} --- it overestimates $\beta_1$.
\end{enumerate}

\vfill
\hfill\hyperlink{q6}{\beamerbutton{Back to Problem 6}}
\end{frame}

% ---- S7 ----
\begin{frame}{Solution 7: Attenuation bias}
\label{s7}
\begin{enumerate}[(a)]
  \item $Y_i = \beta_0 + \beta_1(X_i - w_i) + \varepsilon_i = \beta_0 + \beta_1 X_i + \underbrace{(\varepsilon_i - \beta_1 w_i)}_{u_i}$.

  \item $\operatorname{cov}(X, u) = \operatorname{cov}(X^* + w,\; \varepsilon - \beta_1 w) = -\beta_1\sigma_w^2$.

  \item $\widehat{\beta}_1^{OLS} \xrightarrow{p} \beta_1 + \dfrac{-\beta_1\sigma_w^2}{\sigma_{X^*}^2 + \sigma_w^2} = \beta_1\underbrace{\left(\dfrac{\sigma_{X^*}^2}{\sigma_{X^*}^2 + \sigma_w^2}\right)}_{< 1}$.

  The coefficient is ``shrunk'' toward zero --- \alert{attenuation bias}. More noise ($\sigma_w^2 \uparrow$) $\Rightarrow$ more attenuation.

  \item $\text{plim} = 2 \times \dfrac{4}{4+1} = 2 \times \dfrac{4}{5} = \boxed{1.6}$. OLS underestimates by 20\%.
\end{enumerate}

\vfill
\hfill\hyperlink{q7}{\beamerbutton{Back to Problem 7}}
\end{frame}

% ---- S8 ----
\begin{frame}{Solution 8: Simultaneous equations bias (1/2)}
\label{s8}
\begin{enumerate}[(a)]
  \item Set demand $=$ supply: $\alpha_0 + \alpha_1 P + u^d = \gamma_0 + \gamma_1 P + u^s$.

  Solve: $P_i = \dfrac{(\gamma_0 - \alpha_0) + (u_i^s - u_i^d)}{\alpha_1 - \gamma_1}$.

  \item $\operatorname{cov}(P, u^d) = \dfrac{1}{\alpha_1 - \gamma_1}[\operatorname{cov}(u^s, u^d) - \operatorname{var}(u^d)] = \dfrac{-\sigma_d^2}{\alpha_1 - \gamma_1}$.

  Since $\alpha_1 < 0 < \gamma_1$: $\alpha_1 - \gamma_1 < 0$, so $\operatorname{cov}(P, u^d) > 0$.

  \emph{Intuition}: A positive demand shock raises the equilibrium price.
\end{enumerate}

\vfill
\hfill\hyperlink{s8b}{\beamerbutton{Continued}}
\hfill\hyperlink{q8}{\beamerbutton{Back to Problem 8}}
\end{frame}

\begin{frame}{Solution 8: Simultaneous equations bias (2/2)}
\label{s8b}
\begin{enumerate}[(a)]
  \setcounter{enumi}{2}
  \item $\operatorname{var}(P) = \dfrac{\sigma_s^2 + \sigma_d^2}{(\alpha_1 - \gamma_1)^2}$.

  \[
  \widehat{\alpha}_1^{OLS} \xrightarrow{p} \alpha_1 + \frac{\operatorname{cov}(P,u^d)}{\operatorname{var}(P)} = \alpha_1 + \frac{-\sigma_d^2(\alpha_1 - \gamma_1)}{(\sigma_s^2+\sigma_d^2)} = \frac{\alpha_1\sigma_s^2 + \gamma_1\sigma_d^2}{\sigma_s^2 + \sigma_d^2}
  \]

  \item This is a \alert{variance-weighted average} of the demand slope $\alpha_1$ and the supply slope $\gamma_1$. OLS estimates \emph{neither} the demand nor the supply curve --- it produces a meaningless hybrid. Only when $\sigma_d^2 = 0$ (no demand shocks) does OLS recover $\alpha_1$.
\end{enumerate}

\vfill
\hfill\hyperlink{q8}{\beamerbutton{Back to Problem 8}}
\end{frame}

% ---- S9 ----
\begin{frame}{Solution 9: Exogeneity failure and TSLS inconsistency}
\label{s9}
\begin{enumerate}[(a)]
  \item From $\operatorname{cov}(Z, Y) = \beta_1\operatorname{cov}(Z,X) + \operatorname{cov}(Z,u)$:
  \[
  \frac{\operatorname{cov}(Z,Y)}{\operatorname{cov}(Z,X)} = \beta_1 + \frac{\rho}{\operatorname{cov}(Z,X)}
  \]
  Since $s_{ZY}/s_{ZX} \xrightarrow{p}$ this ratio: $\widehat{\beta}_1^{TSLS} \xrightarrow{p} \beta_1 + \dfrac{\rho}{\operatorname{cov}(Z,X)}$. \checkmark

  \item Bias $= 0.01/0.5 = \boxed{0.02}$. Small.

  \item Bias $= 0.01/0.02 = \boxed{0.50}$. Twenty-five times larger!

  \item The asymptotic bias is $\rho/\operatorname{cov}(Z,X)$. When the instrument is weak ($\operatorname{cov}(Z,X) \approx 0$), the denominator shrinks, \alert{amplifying} even tiny exogeneity violations into large biases.
\end{enumerate}

\vfill
\hfill\hyperlink{q9}{\beamerbutton{Back to Problem 9}}
\end{frame}

% ---- S10 ----
\begin{frame}{Solution 10: Direction of OVB with IV}
\label{s10}
\begin{enumerate}[(a)]
  \item $\operatorname{cov}(Z, u) = \beta_2\operatorname{cov}(Z, W) + \operatorname{cov}(Z, \varepsilon) = 0$. Since $\operatorname{cov}(Z,\varepsilon) = 0$ by assumption, we need $\beta_2\operatorname{cov}(Z, W) = 0$.

  \item Since $\beta_2 \neq 0$ (otherwise there is no OVB), we must have $\operatorname{cov}(Z, W) = 0$. The instrument must be \alert{uncorrelated with the omitted variable}.

  \item \textbf{Parental income is a bad instrument.} It is correlated with ability ($W$): children from wealthier families tend to have advantages (better nutrition, tutoring, etc.)\ that correlate with ability. So $\operatorname{cov}(Z, W) \neq 0$, violating exogeneity. A better instrument: quarter of birth (Angrist \& Krueger, 1991).
\end{enumerate}

\vfill
\hfill\hyperlink{q10}{\beamerbutton{Back to Problem 10}}
\end{frame}

% ---- S11 ----
\begin{frame}{Solution 11: The TSLS formula}
\label{s11}
\begin{enumerate}[(a)]
  \item Since $\widehat{X}_i = \widehat{\pi}_0 + \widehat{\pi}_1 Z_i$, deviations are $\widehat{X}_i - \bar{\widehat{X}} = \widehat{\pi}_1(Z_i - \bar{Z})$. Substitute:
  \begin{align*}
  s_{\widehat{X}Y} &= \frac{1}{n-1}\sum \widehat{\pi}_1(Z_i - \bar{Z})(Y_i - \bar{Y}) = \widehat{\pi}_1 \, s_{ZY} \\
  s_{\widehat{X}}^2 &= \frac{1}{n-1}\sum \widehat{\pi}_1^2(Z_i - \bar{Z})^2 = \widehat{\pi}_1^2 \, s_Z^2
  \end{align*}
  $\widehat{\beta}_1^{TSLS} = \dfrac{\widehat{\pi}_1 s_{ZY}}{\widehat{\pi}_1^2 s_Z^2} = \dfrac{s_{ZY}}{\widehat{\pi}_1 s_Z^2} = \dfrac{s_{ZY}}{s_{ZX}}$, \quad since $\widehat{\pi}_1 = s_{ZX}/s_Z^2$. \checkmark

  \item \textbf{Wald estimator}: $s_{ZY}$ is the \emph{reduced-form} effect ($Z \to Y$) and $s_{ZX}$ is the \emph{first-stage} effect ($Z \to X$). The ratio recovers the structural effect $X \to Y$: ``how much $Y$ moves per unit of $X$ movement caused by $Z$.''
\end{enumerate}

\vfill
\hfill\hyperlink{q11}{\beamerbutton{Back to Problem 11}}
\end{frame}

% ---- S12 ----
\begin{frame}{Solution 12: Consistency of TSLS}
\label{s12}
\begin{enumerate}[(a)]
  \item $\operatorname{cov}(Z, Y) = \operatorname{cov}(Z, \beta_0 + \beta_1 X + u) = \beta_1\operatorname{cov}(Z,X) + \underbrace{\operatorname{cov}(Z,u)}_{=0 \text{ (exogeneity)}}$.

  Solve: $\beta_1 = \dfrac{\operatorname{cov}(Z,Y)}{\operatorname{cov}(Z,X)}$ \quad (requires $\operatorname{cov}(Z,X) \neq 0$: \alert{relevance}).

  \item By LLN: $s_{ZY} \xrightarrow{p} \operatorname{cov}(Z,Y)$, $s_{ZX} \xrightarrow{p} \operatorname{cov}(Z,X)$.

  By continuous mapping theorem (Slutsky): $\dfrac{s_{ZY}}{s_{ZX}} \xrightarrow{p} \dfrac{\operatorname{cov}(Z,Y)}{\operatorname{cov}(Z,X)} = \beta_1$. \checkmark

  \item \textbf{Exogeneity}: used in (a) to set $\operatorname{cov}(Z,u) = 0$, allowing isolation of $\beta_1$.

  \textbf{Relevance}: used in (a) to divide by $\operatorname{cov}(Z,X) \neq 0$, and in (b) to apply Slutsky (denominator converges to a nonzero constant, not zero).
\end{enumerate}

\vfill
\hfill\hyperlink{q12}{\beamerbutton{Back to Problem 12}}
\end{frame}

% ---- S13 ----
\begin{frame}{Solution 13: First-stage $F$-statistic}
\label{s13}
\begin{enumerate}[(a)]
  \item Yes, $F = 7.2 < 10$. The instrument is \alert{weak} by the Stock--Yogo rule of thumb.

  \item Fraction of OLS bias $\approx \dfrac{1}{E(F) - 1} = \dfrac{1}{7.2 - 1} = \dfrac{1}{6.2} \approx \boxed{16.1\%}$. TSLS retains about 16\% of the OLS bias --- substantial.

  \item Three strategies:
  \begin{itemize}
    \item Find stronger instruments (requires domain knowledge).
    \item If overidentified, discard the weakest instruments.
    \item Use weak-instrument-robust inference: Anderson--Rubin or CLR tests/confidence sets.
  \end{itemize}

  \item With weak instruments, the TSLS distribution is \alert{not well-approximated by the normal}. It is heavy-tailed and centered closer to the OLS plim, so standard $t$-tests and CIs have incorrect size (true coverage $\ll$ nominal 95\%).
\end{enumerate}

\vfill
\hfill\hyperlink{q13}{\beamerbutton{Back to Problem 13}}
\end{frame}

% ---- S14 ----
\begin{frame}{Solution 14: Identification and the order condition}
\label{s14}
\begin{enumerate}[(a)]
  \item $k = 2$ endogenous regressors $\Rightarrow$ need $m \geq 2$ instruments. \textbf{Minimum: 2 instruments.}

  \item $m = 3 > k = 2$: the model is \alert{overidentified} (one extra instrument).

  \item Two first-stage regressions (one per endogenous variable):
  \begin{align*}
  X_{1i} &= \pi_0 + \pi_1 Z_{1i} + \pi_2 Z_{2i} + \pi_3 Z_{3i} + \pi_4 W_{1i} + v_{1i} \\
  X_{2i} &= \delta_0 + \delta_1 Z_{1i} + \delta_2 Z_{2i} + \delta_3 Z_{3i} + \delta_4 W_{1i} + v_{2i}
  \end{align*}
  Each includes \alert{all} instruments and \alert{all} exogenous regressors.

  \item Yes, because $m > k$. The $J$-test has $m - k = 3 - 2 = \boxed{1}$ degree of freedom. Under $H_0$: $J \sim \chi^2_1$.
\end{enumerate}

\vfill
\hfill\hyperlink{q14}{\beamerbutton{Back to Problem 14}}
\end{frame}

% ---- S15 ----
\begin{frame}{Solution 15: The $J$-test}
\label{s15}
\begin{enumerate}[(a)]
  \item Steps:
  \begin{enumerate}[1.]
    \item Estimate the model by TSLS; obtain residuals $\widehat{u}_i^{TSLS}$.
    \item Regress $\widehat{u}_i^{TSLS}$ on all $Z$'s and $W$'s. Compute the $F$-statistic testing all $Z$-coefficients $= 0$.
    \item $J = m \times F$.
  \end{enumerate}

  \item $J \sim \chi^2_{m-k} = \chi^2_{3-1} = \chi^2_2$.

  \item $J = 7.8 > 5.99 = \chi^2_2(0.05)$. \alert{Reject $H_0$} at the 5\% level.

  \item No. The $J$-test rejects the hypothesis that \emph{all} instruments are exogenous, but it could be that only \emph{one} is endogenous. It does not identify \emph{which} instrument(s) are problematic. Moreover, the test is conditional on at least $k$ instruments being valid --- it can never confirm the baseline identifying assumption.
\end{enumerate}

\vfill
\hfill\hyperlink{q15}{\beamerbutton{Back to Problem 15}}
\end{frame}

% ---- S16 ----
\begin{frame}{Solution 16: Sampling distribution of TSLS}
\label{s16}
\begin{enumerate}[(a)]
  \item $E(\bar{q}) \approx E\!\left[\frac{1}{n}\sum(Z_i - \mu_Z)u_i\right] = E[(Z_i - \mu_Z)u_i] = \operatorname{cov}(Z,u) = 0$ under exogeneity.

  Since $q_i = (Z_i - \mu_Z)u_i$ are i.i.d.\ with $E(q_i) = 0$ and $\operatorname{var}(q_i) = \sigma_q^2 < \infty$, the \alert{Central Limit Theorem} gives $\sqrt{n}\,\bar{q} \xrightarrow{d} N(0, \sigma_q^2)$.

  \item \textbf{Slutsky's theorem}: since the denominator $\xrightarrow{p} \operatorname{cov}(Z,X) \neq 0$ (a constant), the ratio of an asymptotically normal numerator and a probability limit in the denominator is asymptotically normal.

  \item $\sigma_{\widehat{\beta}_1^{TSLS}}^2 = \dfrac{1}{n} \cdot \dfrac{\operatorname{var}[(Z_i - \mu_Z)u_i]}{[\operatorname{cov}(Z,X)]^2}$.

  When $\operatorname{cov}(Z,X)$ is small, the denominator shrinks, inflating the variance. Weak instruments $\Rightarrow$ \alert{imprecise} estimates.
\end{enumerate}

\vfill
\hfill\hyperlink{q16}{\beamerbutton{Back to Problem 16}}
\end{frame}

% ---- S17 ----
\begin{frame}{Solution 17: Cigarette demand application}
\label{s17}
\begin{enumerate}[(a)]
  \item All $F$-statistics $\gg 10$ (33.7, 107.2, 88.6). Instruments are \alert{not weak}.

  \item $J$-test $p = 0.026 < 0.05$: reject $H_0$ that both instruments are exogenous. At least one instrument is likely endogenous.

  \item \textbf{Prefer SalesTax} ($\widehat{\beta}_1 = -0.94$). General sales tax is set for broad public finance reasons, unrelated to the cigarette market. Cigarette-specific tax may be influenced by tobacco industry lobbying and is thus more likely correlated with demand-side factors.

  \item A 1\% price increase reduces cigarette consumption by about 0.94\% in the long run. Demand is approximately unit-elastic at the 10-year horizon. To reduce consumption by 20\%, prices need to rise by $\approx 20/0.94 \approx 21\%$.
\end{enumerate}

\vfill
\hfill\hyperlink{q17}{\beamerbutton{Back to Problem 17}}
\end{frame}

% ---- S18 ----
\begin{frame}{Solution 18: Anderson--Rubin test}
\label{s18}
\begin{enumerate}[(a)]
  \item If $\beta_1 = \beta_{1,0}$, then $Y_i^* = Y_i - \beta_{1,0}X_i = \beta_0 + u_i + \beta_2 W_{1i} + \cdots$ Under exogeneity, $\operatorname{cov}(Z, u) = 0$. So $Z$ has no predictive power for $Y^*$ beyond $W$'s --- the $F$-statistic should be small.

  \item Standard TSLS inference relies on the denominator converging to a nonzero constant (Slutsky). With weak instruments, this fails. The AR test uses the $F$-distribution of a regression, which is valid \alert{regardless} of instrument strength --- it never divides by $\operatorname{cov}(Z,X)$.

  \item For each candidate $\beta_{1,0}$, run the AR test. The 95\% confidence set is $\{\beta_{1,0} : \text{AR test does not reject at 5\%}\}$.

  \item AR confidence sets can be: \emph{empty} (no $\beta_{1,0}$ is consistent with the data), \emph{unbounded} (instruments are too weak to pin down $\beta_1$), or \emph{disconnected} (two disjoint intervals).
\end{enumerate}

\vfill
\hfill\hyperlink{q18}{\beamerbutton{Back to Problem 18}}
\end{frame}

% ---- S19 ----
\begin{frame}{Solution 19: IV with simultaneous causality (1/2)}
\label{s19}
\begin{enumerate}[(a)]
  \item Set demand $=$ supply:
  \[
  \alpha_0 + \alpha_1 P + u^d = \gamma_0 + \gamma_1 P + \delta Z + u^s
  \]
  Solve: $\boxed{P_i = \dfrac{(\gamma_0 - \alpha_0) + \delta Z_i + (u_i^s - u_i^d)}{\alpha_1 - \gamma_1}}$.

  \item $\operatorname{cov}(Z, P) = \dfrac{\delta\operatorname{var}(Z) + \operatorname{cov}(Z, u^s) - \operatorname{cov}(Z, u^d)}{\alpha_1 - \gamma_1} = \dfrac{\delta\operatorname{var}(Z)}{\alpha_1 - \gamma_1} \neq 0$

  since $\delta \neq 0$ and $\operatorname{var}(Z) > 0$. \alert{Relevance holds.} \checkmark

  \item $\operatorname{cov}(Z, u^d) = 0$ by assumption. \alert{Exogeneity holds.} \checkmark
\end{enumerate}

\vfill
\hfill\hyperlink{s19b}{\beamerbutton{Continued}}
\hfill\hyperlink{q19}{\beamerbutton{Back to Problem 19}}
\end{frame}

\begin{frame}{Solution 19: IV with simultaneous causality (2/2)}
\label{s19b}
\begin{enumerate}[(a)]
  \setcounter{enumi}{3}
  \item From the demand equation: $Q_i = \alpha_0 + \alpha_1 P_i + u_i^d$.

  Take covariances with $Z$:
  \[
  \operatorname{cov}(Z, Q) = \alpha_1\operatorname{cov}(Z, P) + \underbrace{\operatorname{cov}(Z, u^d)}_{=0}
  \]

  Therefore:
  \[
  \frac{\operatorname{cov}(Z, Q)}{\operatorname{cov}(Z, P)} = \alpha_1 \qquad \checkmark
  \]

  TSLS consistently estimates the demand slope $\alpha_1$ by using only the supply-driven variation in prices (caused by $Z$). The demand shocks $u^d$ are ``projected out.''
\end{enumerate}

\vfill
\hfill\hyperlink{q19}{\beamerbutton{Back to Problem 19}}
\end{frame}

% ---- S20 ----
\begin{frame}{Solution 20: Comprehensive analysis (1/2)}
\label{s20}
\begin{enumerate}[(a)]
  \item OLS is biased \alert{downward}. Sicker patients receive more spending ($\operatorname{cov}(X, u) < 0$ because sicker patients have worse outcomes, i.e., $u$ captures unobserved health). This is a combination of OVB (unobserved severity) and simultaneous causality (spending responds to health status). OLS shows spending \emph{hurts} ($-0.05$), but this reflects selection, not causation.

  \item $F = 14.3 > 10$. Instruments pass the weak-instrument test. TSLS standard inference is reliable.

  \item $J = 0.45$, $p = 0.50$: \alert{do not reject} $H_0$ that both instruments are exogenous. This is \emph{consistent with} validity but does not \emph{prove} it --- the $J$-test can only detect inconsistencies between instruments, not absolute violations.
\end{enumerate}

\vfill
\hfill\hyperlink{s20b}{\beamerbutton{Continued}}
\hfill\hyperlink{q20}{\beamerbutton{Back to Problem 20}}
\end{frame}

\begin{frame}{Solution 20: Comprehensive analysis (2/2)}
\label{s20b}
\begin{enumerate}[(a)]
  \setcounter{enumi}{3}
  \item \textbf{Distance to teaching hospital} ($Z_1$): Plausibly exogenous --- residential location is largely determined before a heart attack occurs. Relevance: patients closer to teaching hospitals receive more intensive (expensive) care. Concern: residential sorting (healthier/wealthier people may live near hospitals).

  \textbf{Property tax revenue} ($Z_2$): Relevant (higher tax revenue $\to$ better-funded hospitals $\to$ more spending). But potentially endogenous: wealthy areas have higher property taxes \emph{and} healthier populations. \alert{$Z_1$ is more credible.}

  \item OLS: $\widehat{\beta}_1 = -0.05$ (spending appears harmful --- selection bias dominates). TSLS: $\widehat{\beta}_1 = 0.12$ (spending improves survival). The sign reversal is exactly what we expect when the dominant endogeneity is negative (sicker patients get more spending). After removing this contamination, IV reveals the \alert{true positive causal effect} of hospital spending on survival.
\end{enumerate}

\vfill
\hfill\hyperlink{q20}{\beamerbutton{Back to Problem 20}}
\end{frame}

\end{document}
